\section{Stabilne $\infty$-kategorie i wyższa algebra homologiczna}\label{cat:sec:stable-infty-categories}

W kategorii triangulowanej (rozdział~\ref{cat:sec:triangulated-derived}) przesunięcia i trójkąty wyróżnione opisują konstrukcje homologiczne na poziomie klas homotopii. Nie rejestrują jednak wszystkich homotopii między morfizmami ani wyższych zależności między nimi. Stabilne $\infty$-kategorie zachowują tę informację i pozwalają definiować włókna oraz kowłókna przez własności uniwersalne na poziomie przestrzeni odwzorowań.

Pracujemy w modelu quasi-kategorii z rozdziału~\ref{cat:sec:quasicategories}, w ustalonych uniwersach. Przez $\mathcal S$ oznaczamy $\infty$-kategorię przestrzeni (typów homotopii, reprezentowanych przez kompleksy Kana), a przez $\operatorname{Map}_{\mathcal C}(X,Y)$ przestrzeń odwzorowań. Zapis $\simeq$ oznacza równoważność obiektów lub przestrzeni, nie ich dosłowną równość. Diagram przemienny w $\infty$-kategorii zawiera również dane odpowiednich homotopii: kwadrat jest diagramem $\Delta^1\times\Delta^1\to\mathcal C$, a nie jedynie kwadratem przemiennym w kategorii homotopii.

\subsection{Kwadraty kartezjańskie i kokartezjańskie}

Przyjmujemy następujące podstawowe wyniki rachunku granic w $\infty$-kategoriach (por. J.~Lurie, \emph{Higher Topos Theory}, rozdział 4, oraz \emph{Kerodon}, podrozdział 7.6.2). Uzupełniają one definicję~\ref{cat:def:quasicat-limits}.

\begin{twierdzenie}[Kryterium granic przez przestrzenie odwzorowań; wynik zewnętrzny]\label{cat:thm:stable-mapping-limit}
Niech $D:K\to\mathcal C$ będzie diagramem. Stożek o wierzchołku $L$ jest stożkiem granicznym wtedy i tylko wtedy, gdy dla każdego $T\in\mathcal C$ indukowany morfizm
\[
 \operatorname{Map}_{\mathcal C}(T,L)\longrightarrow
 \lim_{k\in K}\operatorname{Map}_{\mathcal C}(T,D(k))
\]
jest równoważnością w $\mathcal S$. Dualnie, stożek kograniczny o wierzchołku $M$ jest uniwersalny wtedy i tylko wtedy, gdy dla każdego $T$ indukowany morfizm
\[
 \operatorname{Map}_{\mathcal C}(M,T)\longrightarrow
 \lim_{k\in K^{\mathrm{op}}}\operatorname{Map}_{\mathcal C}(D(k),T)
\]
jest równoważnością.
\end{twierdzenie}

Granice po prawej stronie są granicami w $\mathcal S$, nie granicami zbiorów klas składowych spójnych. W szczególności nie wolno zastąpić iloczynu włóknistego przestrzeni zwykłym iloczynem włóknistym zbiorów.

\begin{definicja}[Kwadrat kartezjański, kokartezjański i bikartezjański]\label{cat:def:stable-bicartesian}
Rozważmy kwadrat
\begin{center}
\begin{tikzpicture}[>=stealth]
 \node (a) at (0,1.8) {$A$};
 \node (b) at (3.4,1.8) {$B$};
 \node (c) at (0,0) {$C$};
 \node (d) at (3.4,0) {$D$};
 \draw[->] (a)--node[above] {$f$}(b);
 \draw[->] (a)--node[left] {$g$}(c);
 \draw[->] (b)--node[right] {$h$}(d);
 \draw[->] (c)--node[below] {$k$}(d);
\end{tikzpicture}
\end{center}
Jest kartezjański, gdy wskazany stożek realizuje $A\simeq B\times_D C$. Jest kokartezjański, gdy wskazany stożek kograniczny realizuje $D\simeq B\sqcup_A C$. Jest bikartezjański, gdy spełnia oba warunki.
\end{definicja}

Przez skończone granice i kogranice rozumiemy granice i kogranice diagramów indeksowanych skończonymi zbiorami symplicjalnymi (mającymi skończenie wiele niezdegenerowanych sympleksów). Nie jest to to samo co żądanie granic po nerwach wszystkich kategorii o skończonej liczbie obiektów i morfizmów: taki nerw może mieć niezdegenerowane sympleksy w dowolnie wysokim wymiarze.

\begin{twierdzenie}[Generowanie skończonych granic i kogranic; wynik zewnętrzny]\label{cat:thm:stable-finite-limit-generation}
$\infty$-kategoria ma skończone granice wtedy i tylko wtedy, gdy ma obiekt końcowy i iloczyny włókniste. Funktor między $\infty$-kategoriami mającymi skończone granice zachowuje je wtedy i tylko wtedy, gdy zachowuje obiekt końcowy i kwadraty kartezjańskie. Dualnie, dla skończonych kogranic wystarczają obiekt początkowy i sumy amalgamowane.
\end{twierdzenie}
Por. \emph{Kerodon}, wniosek 7.6.2.42, oraz dualny wynik dla kogranic; również J.~Lurie, \emph{Higher Topos Theory}, wnioski 4.4.2.4 i 4.4.2.5. Będziemy stosować te wyniki bez powtarzania ich ogólnych dowodów symplicjalnych.

\begin{lemat}[Sklejanie kwadratów]\label{cat:lem:stable-square-pasting}
W diagramie złożonym z dwóch sąsiednich kwadratów kokartezjańskich kwadrat zewnętrzny jest kokartezjański. Jeżeli lewy kwadrat i kwadrat zewnętrzny są kokartezjańskie, to prawy kwadrat jest kokartezjański. Dualnie dla kwadratów kartezjańskich zachodzi reguła sklejania i skracania z prawym kwadratem.
\end{lemat}
\begin{proof}
Dla diagramu o górnym wierszu $A\to B\to C$ i dolnym wierszu $D\to E\to F$ zakładamy, że $E\simeq B\sqcup_A D$ za pomocą lewego kwadratu. Dla dowolnego $T$ oznaczmy $\operatorname{Map}_{\mathcal C}(Z,T)$ przez $M_Z$. Kryterium~\ref{cat:thm:stable-mapping-limit} daje $M_E\simeq M_B\times_{M_A}M_D$. Wówczas
\[
 M_C\times_{M_B}M_E
 \simeq M_C\times_{M_B}(M_B\times_{M_A}M_D)
 \simeq M_C\times_{M_A}M_D.
\]
Druga równoważność wynika z własności uniwersalnej granicy: obie strony opisują punkt w $M_C$, punkt w $M_D$ i ścieżkę między ich obrazami w $M_A$, wraz z wyższymi homotopiami. Równoważność jest zgodna z morfizmami z $M_F$ indukowanymi przez diagram. Zatem morfizm testujący kokartezjańskość prawego kwadratu jest równoważnością dla każdego $T$ wtedy i tylko wtedy, gdy taką własność ma morfizm testujący zewnętrzny kwadrat. Wniosek wynika z kryterium granic przez przestrzenie odwzorowań. Przechodząc do $\mathcal C^{\mathrm{op}}$, dostajemy wersję kartezjańską z odwróconą kolejnością kwadratów.
\end{proof}

\subsection{Obiekt zerowy, włókno i kowłókno}

\begin{definicja}[Punktowana $\infty$-kategoria]\label{cat:def:stable-pointed}
Obiekt zerowy $0$ jest jednocześnie początkowy i końcowy. $\infty$-kategorię mającą obiekt zerowy nazywamy punktowaną. W niej morfizm zerowy $X\to Y$ jest złożeniem $X\to0\to Y$.
\end{definicja}

Przestrzenie $\operatorname{Map}(X,0)$ i $\operatorname{Map}(0,Y)$ są ściągalne. W konsekwencji wybory strzałek w tym złożeniu i ich składania są określone z dokładnością do koherentnej homotopii; dają wyróżniony punkt przestrzeni $\operatorname{Map}(X,Y)$ z dokładnością do ściągalnego wyboru. W kategorii homotopii otrzymujemy dobrze określony morfizm zerowy. W dalszym ciągu przestrzenie odwzorowań traktujemy jako przestrzenie punktowane w ten sposób.

\begin{definicja}[Włókno i kowłókno]\label{cat:def:stable-fiber-cofiber}
Dla $f:X\to Y$ w punktowanej $\infty$-kategorii, jeśli istnieją odpowiednie granice i kogranice, definiujemy
\[
 \operatorname{fib}(f)=X\times_Y0,\qquad
 \operatorname{cofib}(f)=Y\sqcup_X0.
\]
Przez ciąg włóknowy rozumiemy diagram $A\to B\to C$ z określoną homotopią złożenia do morfizmu zerowego, dla którego kwadrat $A\to B$, $0\to C$ jest kartezjański. Przez ciąg kowłóknowy rozumiemy taki diagram, dla którego kwadrat $A\to B$, $0\to C$ jest kokartezjański.
\end{definicja}

W obu przypadkach chodzi o ten sam kształt diagramu:
\begin{center}
\begin{tikzpicture}[>=stealth]
 \node (a) at (0,1.8) {$A$};
 \node (b) at (3.5,1.8) {$B$};
 \node (o) at (0,0) {$0$};
 \node (c) at (3.5,0) {$C$};
 \draw[->] (a)--node[above] {$f$}(b);
 \draw[->] (a)--(o);
 \draw[->] (b)--node[right] {$g$}(c);
 \draw[->] (o)--(c);
\end{tikzpicture}
\end{center}
Zależnie od własności kwadratu mamy $A\simeq\operatorname{fib}(g)$ lub $C\simeq\operatorname{cofib}(f)$. Homotopia zerowa złożenia $gf$ jest częścią danych i nie można jej pominąć.

\begin{stwierdzenie}[Włókna przestrzeni odwzorowań]\label{cat:prop:stable-mapping-fibers}
Dla $f:X\to Y$ i dowolnego $T$ zachodzą naturalne równoważności przestrzeni punktowanych
\[
\begin{aligned}
 \operatorname{Map}(T,\operatorname{fib}(f))
 &\simeq\operatorname{fib}_0\bigl(\operatorname{Map}(T,X)
                    \to\operatorname{Map}(T,Y)\bigr),\\
 \operatorname{Map}(\operatorname{cofib}(f),T)
 &\simeq\operatorname{fib}_0\bigl(\operatorname{Map}(Y,T)
                    \to\operatorname{Map}(X,T)\bigr).
\end{aligned}
\]
Po prawej stronie $\operatorname{fib}_0$ oznacza włókno homotopijne nad wyróżnionym punktem.
\end{stwierdzenie}
\begin{proof}
Z twierdzenia~\ref{cat:thm:stable-mapping-limit} otrzymujemy w pierwszym przypadku
\[
 \operatorname{Map}(T,X\times_Y0)
 \simeq\operatorname{Map}(T,X)
       \times_{\operatorname{Map}(T,Y)}\operatorname{Map}(T,0).
\]
Ostatni czynnik jest ściągalny i jego morfizm do $\operatorname{Map}(T,Y)$ wybiera morfizm zerowy. Jest to właśnie włókno homotopijne. Dla drugiego równania stosujemy część dualną twierdzenia i ściągalność $\operatorname{Map}(0,T)$. Wyróżnione punkty obu stron odpowiadają sobie przez funktorialność konstrukcji.
\end{proof}

\subsection{Stabilność i przesunięcie}

\begin{definicja}[Stabilna $\infty$-kategoria]\label{cat:def:stable-infty-category}
Punktowaną $\infty$-kategorię $\mathcal C$ nazywamy stabilną, jeżeli ma skończone granice i kogranice oraz kwadrat w $\mathcal C$ jest kartezjański wtedy i tylko wtedy, gdy jest kokartezjański.
\end{definicja}
Jest to równoważna postać definicji J.~Lurie, \emph{Higher Algebra}, definicja 1.1.1.9 i stwierdzenie 1.1.3.4. W szczególności ciągi włóknowe i kowłóknowe są w kategorii stabilnej tymi samymi diagramami. Nazywamy je ciągami dokładnymi, mając na myśli dokładność w sensie $\infty$-kategorii.

\begin{definicja}[Zawieszenie i obiekt pętli]\label{cat:def:stable-suspension-loop}
W punktowanej $\infty$-kategorii mającej skończone granice i kogranice kładziemy
\[
 \Sigma X=0\sqcup_X0,\qquad
 \Omega X=0\times_X0.
\]
Są to odpowiednio zawieszenie i obiekt pętli.
\end{definicja}
Konstrukcje są funktorialne z dokładnością do koherentnego wyboru granic i kogranic. Korzystamy tu ze standardowej funktorialności granic: funktor granicy jest prawym sprzężeniem do funktora stałych diagramów, a funktor kogranicy jego lewym odpowiednikiem. Własności uniwersalne zapewniają, że zastosowanie konstrukcji do morfizmów diagramów wyznacza całe funktory $\Sigma$ i $\Omega$, a nie tylko funkcje na klasach obiektów.

\begin{stwierdzenie}[Sprzężenie zawieszenia i pętli]\label{cat:prop:stable-suspension-loop-adjunction}
W punktowanej $\infty$-kategorii mającej skończone granice i kogranice zachodzi $\Sigma\dashv\Omega$.
\end{stwierdzenie}
\begin{proof}
Oznaczmy przez $\Omega_{\mathcal S}$ funktor pętli w przestrzeniach punktowanych. Stosując kryterium~\ref{cat:thm:stable-mapping-limit} dwukrotnie, otrzymujemy naturalne równoważności
\[
\begin{aligned}
 \operatorname{Map}(\Sigma X,Y)
 &\simeq *\times_{\operatorname{Map}(X,Y)}*
 \simeq\Omega_{\mathcal S}\operatorname{Map}(X,Y),\\
 \operatorname{Map}(X,\Omega Y)
 &\simeq *\times_{\operatorname{Map}(X,Y)}*
 \simeq\Omega_{\mathcal S}\operatorname{Map}(X,Y).
\end{aligned}
\]
W obu wzorach morfizmy z przestrzeni jednopunktowej wybierają morfizm zerowy. Złożenie pierwszej równoważności z odwrotnością drugiej daje naturalną równoważność przestrzeni odwzorowań, czyli sprzężenie w sensie rozdziału~\ref{cat:sec:quasicategories}.
\end{proof}

\begin{twierdzenie}[Odwracalność zawieszenia]\label{cat:thm:stable-suspension-equivalence}
W stabilnej $\infty$-kategorii zawieszenie i obiekt pętli są wzajemnie odwrotnymi równoważnościami:
\[
 \Omega\Sigma\simeq\id_{\mathcal C},\qquad
 \Sigma\Omega\simeq\id_{\mathcal C}.
\]
\end{twierdzenie}
\begin{proof}
Z definicji $\Sigma X$ kwadrat po lewej stronie poniższego diagramu jest kokartezjański; z definicji $\Omega Y$ kwadrat po prawej stronie jest kartezjański.
\begin{center}
\begin{tikzpicture}[>=stealth]
 \node (a) at (0,1.8) {$X$};
 \node (b) at (2.4,1.8) {$0$};
 \node (c) at (0,0) {$0$};
 \node (d) at (2.4,0) {$\Sigma X$};
 \draw[->] (a)--(b); \draw[->] (a)--(c);
 \draw[->] (b)--(d); \draw[->] (c)--(d);
 \node (e) at (5,1.8) {$\Omega Y$};
 \node (f) at (7.4,1.8) {$0$};
 \node (g) at (5,0) {$0$};
 \node (h) at (7.4,0) {$Y$};
 \draw[->] (e)--(f); \draw[->] (e)--(g);
 \draw[->] (f)--(h); \draw[->] (g)--(h);
\end{tikzpicture}
\end{center}
Stabilność czyni oba kwadraty bikartezjańskimi. Pierwszy realizuje zatem $X$ jako iloczyn włóknisty $0\times_{\Sigma X}0$, a drugi realizuje $Y$ jako sumę amalgamowaną $0\sqcup_{\Omega Y}0$. Morfizmy porównawcze $X\to\Omega\Sigma X$ i $\Sigma\Omega Y\to Y$ są równoważnościami, ponieważ porównują dwa stożki graniczne lub dwa stożki kograniczne. Są to jednostka i kojednostka sprzężenia ze stwierdzenia~\ref{cat:prop:stable-suspension-loop-adjunction}: pierwszy morfizm odpowiada identyczności $\Sigma X\to\Sigma X$, a drugi identyczności $\Omega Y\to\Omega Y$, przez podane tam własności uniwersalne. Naturalność jednostki i kojednostki daje równoważności funktorów.
\end{proof}

\begin{przyklad}[Punktowanie nie wystarcza do stabilności]\label{cat:ex:stable-pointed-counterexample}
Zwykła kategoria grup abelowych ma obiekt zerowy oraz wszystkie skończone granice i kogranice. W jej nerwie granice i kogranice diagramów zwykłych odpowiadają granicom i kogranicom w kategorii grup abelowych (przykład~\ref{cat:ex:quasicat-ordinary-limits}). Dla dowolnej grupy $A$ suma amalgamowana $0\sqcup_A0$ jest zerowa. Istotnie, dla każdego $B$ istnieje dokładnie jedna zgodna para morfizmów $0\to B$. Jeżeli $A$ jest niezerowa, to kwadrat $A\to0$, $0\to0$ nie jest kartezjański, gdyż iloczyn włóknisty $0\times_00$ jest zerowy. Nerw kategorii grup abelowych nie jest więc stabilną $\infty$-kategorią. Wyjaśnia to, dlaczego w wyższej algebrze homologicznej stosujemy konstrukcje pochodne, a nie zwykłe iloczyny włókniste i sumy amalgamowane grup abelowych.
\end{przyklad}

\subsection{Funktory dokładne}

\begin{definicja}[Dokładność funktora stabilnego]\label{cat:def:stable-exact-functor}
Funktor pomiędzy stabilnymi $\infty$-kategoriami nazywamy dokładnym, jeżeli zachowuje skończone granice i skończone kogranice.
\end{definicja}

\begin{stwierdzenie}[Kryteria dokładności]\label{cat:prop:stable-exact-criteria}
Dla funktora $F:\mathcal C\to\mathcal D$ pomiędzy stabilnymi $\infty$-kategoriami równoważne są warunki:
\begin{enumerate}
 \item $F$ jest dokładny;
 \item $F$ zachowuje skończone granice;
 \item $F$ zachowuje skończone kogranice;
 \item $F$ zachowuje obiekt zerowy oraz kwadraty bikartezjańskie.
\end{enumerate}
\end{stwierdzenie}
\begin{proof}
Warunek pierwszy implikuje drugi i trzeci. Jeśli funktor zachowuje skończone granice, to zachowuje obiekt końcowy, będący w obu kategoriach obiektem zerowym. Zachowuje też kwadraty kartezjańskie, a przez stabilność są to dokładnie kwadraty bikartezjańskie. Zatem drugi warunek implikuje czwarty. Analogicznie trzeci warunek implikuje czwarty, wykorzystując obiekt początkowy i kwadraty kokartezjańskie. Jeżeli zachodzi czwarty warunek, to funktor zachowuje obiekt początkowy i końcowy oraz kwadraty kartezjańskie i kokartezjańskie, ponieważ w obu kategoriach oba rodzaje kwadratów się pokrywają. Twierdzenie~\ref{cat:thm:stable-finite-limit-generation} daje zachowywanie wszystkich skończonych granic i kogranic.
\end{proof}

\begin{wniosek}[Sprzężenia stabilne]\label{cat:cor:stable-adjoint-exact}
Jeżeli $F:\mathcal C\rightleftarrows\mathcal D:G$ jest sprzężeniem pomiędzy stabilnymi $\infty$-kategoriami, to zarówno $F$, jak i $G$ są dokładne.
\end{wniosek}
\begin{proof}
Lewe sprzężenie zachowuje istniejące kogranice, a prawe istniejące granice. Dla przykładu dla diagramu $D:K\to\mathcal C$ i jego kogranicy $M$ mamy
\[
 \operatorname{Map}_{\mathcal D}(FM,T)
 \simeq\operatorname{Map}_{\mathcal C}(M,GT)
 \simeq\lim_{k\in K^{\mathrm{op}}}
           \operatorname{Map}_{\mathcal C}(D(k),GT)
 \simeq\lim_{k\in K^{\mathrm{op}}}
           \operatorname{Map}_{\mathcal D}(FD(k),T).
\]
Równoważności są naturalne i zgodne z morfizmami stożka kogranicznego, więc kryterium~\ref{cat:thm:stable-mapping-limit} dowodzi zachowywania kogranicy. Argument dualny daje zachowywanie granic przez $G$. Stwierdzenie~\ref{cat:prop:stable-exact-criteria} kończy dowód.
\end{proof}
